% !TEX program = xelatex
\documentclass[a4paper,11pt]{article}
\usepackage[margin=2cm]{geometry}
\usepackage{amsmath,amssymb,graphicx,caption,float,microtype}
\usepackage[hidelinks]{hyperref}
\graphicspath{{../../random_ising/results/run_004/figures/}}
\captionsetup{font=small,labelfont=bf}
\setlength{\parskip}{0.35em}
\title{Monte Carlo Results for the Two-Dimensional\\Random-Bond Ising Model}
\author{Zihao Xu}
\date{September 17, 2026}
\begin{document}
\maketitle
\vspace{-2em}
\section{Model and the Griffiths region}
We consider spins $s_i=\pm1$ on an $L\times L$ square lattice with periodic boundary conditions and
$N=L^2$ sites. Setting $k_{\mathrm B}=1$, the zero-field Hamiltonian is
\begin{equation}
H=-\sum_{\langle ij\rangle}J_{ij}s_i s_j,\qquad J_{ij}=\begin{cases}1,&\text{with probability }p,\\0,&\text{with probability }1-p.\end{cases}
\end{equation}
Bonds are drawn independently and remain fixed during thermal sampling. This is the bond-diluted
limit of the ferromagnetic random-bond Ising model. Dataset \texttt{run\_004} uses $p=0.5$. We
denote thermal averages in realization $a$ by $\langle\cdot\rangle_a$ and disorder averages by
$[\cdot]_{\rm dis}$.

Disorder allows unusually well-connected, nearly pure regions inside a globally paramagnetic system.
In the Griffiths region, $T_c(p)<T<T_G$, these rare regions can generate enhanced responses and slow
relaxation, while the equilibrium free energy is nonanalytic in magnetic field at zero field. Here
$T_G=2/\ln(1+\sqrt2)\simeq2.269$ is the pure square-lattice transition temperature. Classical
Griffiths singularities can be weak: slow dynamics or broad response distributions alone do not
establish their asymptotic form.

The present concentration equals the square-lattice bond-percolation threshold, $p_c=1/2$, where the
thermodynamic ferromagnetic transition temperature is zero. A heat-capacity maximum near $T=1$
therefore does not imply finite-temperature ordering. Critical connectivity, rare dense regions, and
finite observation times must all be considered in interpreting the data.

\begin{figure}[H]
\centering\includegraphics[width=0.56\linewidth]{phase_diagram.png}
\caption{Schematic temperature--bond-concentration phase diagram. The green curve denotes $T_c(p)$,
separating ferromagnetic order from the paramagnetic phase. The red line is the pure-system
transition temperature, labeled $T_f$ in the image and denoted $T_G$ here. The lightly shaded region
between these boundaries is the Griffiths region. The label $p_{\rm th.}$ corresponds to the
percolation threshold $p_c=1/2$. Black markers illustrate a temperature scan at $p=0.6$; they are
not the $p=0.5$ data reported below.}\label{fig:phase}
\end{figure}
Figure~\ref{fig:phase} emphasizes that the Griffiths region is globally paramagnetic despite its
slowly relaxing, locally correlated rare regions. At fixed $p>p_c$, cooling crosses $T_G$ before
reaching the ordering temperature $T_c(p)$. For the present $p=p_c$, the latter boundary reaches
zero, leaving $0<T<T_G$ without bulk ferromagnetic order. The green curve is schematic and is not a
fit to this dataset.

\section{Monte Carlo procedure}
The simulations used SpinMonteCarlo.jl 1.2.2. Each realization was thermalized with $2048$
Swendsen--Wang (SW) updates: parallel neighboring spins are linked with probability
$1-e^{-2J_{ij}/T}$, and each cluster is independently assigned either orientation. Measurements then
used $65536$ random-site heat-bath sweeps. Each sweep selects $N$ sites with replacement and
resamples them according to
\begin{equation}
P(s_i=+1\mid\{s_j\})=\frac{1+\tanh(h_i/T)}{2},\qquad h_i=\sum_{j:\langle ij\rangle}J_{ij}s_j.
\end{equation}
Local responses also used SW thermalization followed by heat-bath measurements. This is the
implementation that generated this dataset, before the subsequent change to SW sampling of static
observables.

The sizes are $L=10,16,24,32,40,50$, with $T=0.01,0.02,\ldots,3.00$ and $1000$ disorder realizations
per size. The same bond realizations are reused across temperatures. Measurements are divided into
$256$ blocks of $256$ sweeps for block statistics and jackknife estimates. Disorder-average bands
use the sample standard deviation divided by $\sqrt{1000}$; this scatter already includes MC noise.
A fixed block length need not resolve low-temperature slow modes and does not guarantee converged
errors.

\clearpage
\section{Results}
\subsection{Disorder-averaged heat capacity and susceptibility}
For $M=\sum_i s_i$, the observables are
\begin{equation}
C_a=\frac{\langle H^2\rangle_a-\langle H\rangle_a^2}{T^2},\qquad \chi_a=\frac{\langle M^2\rangle_a}{T}.
\end{equation}
Both are total-system quantities; division by $N$ gives per-site values. The susceptibility uses the
exact symmetry $\langle M\rangle_a=0$ of the finite-volume zero-field ensemble, without subtracting
$\langle|M|\rangle_a^2$.
\begin{figure}[H]
\centering\includegraphics[width=\linewidth]{01_mean_observables.png}
\caption{Disorder-averaged total heat capacity (left) and total susceptibility (right) for six
sizes. Shading denotes $\pm1$ standard error of the disorder mean (SEM). The susceptibility axis is
logarithmic.}\label{fig:mean}
\end{figure}
The heat capacity has a broad maximum at $T=0.99$--$1.03$, with little systematic peak displacement
over these sizes. For $L=50$, the maximum occurs at $T=1.02$, where $[C]_{\rm dis}/N\simeq0.420$.
The per-site peak decreases from approximately $0.435$ at $L=10$ to $0.420$ at $L=50$. The growing
total peak therefore mainly reflects increasing volume, rather than a growing per-site singularity.
Its broad shape is consistent with a crossover in local bond-energy fluctuations. At low
temperature, occupied bonds are mostly satisfied and energy fluctuations are suppressed; at high
temperature the factor $T^{-2}$ suppresses the heat capacity.

The susceptibility rises strongly upon cooling, without an isolated maximum in the displayed range.
A useful low-temperature equilibrium picture is a set of nearly rigid disconnected clusters. A
cluster of $n_c$ spins contributes approximately $n_c^2$ to $\langle M^2\rangle$, giving $\chi\simeq
T^{-1}\sum_c n_c^2$ when cluster orientations are independently sampled. Large clusters have
disproportionate weight, making magnetic response particularly sensitive to connectivity. This is a
physical interpretation, rather than a quantitative fit to the trajectories.

At high temperature, separation between sizes also contains the ordinary extensive factor:
asymptotically $\chi\sim N/T$ for weakly correlated spins. Separation of total-susceptibility curves
alone is consequently insufficient evidence for anomalous scaling. At low temperature, heat-bath
trajectories may fail to average over relative cluster orientations, biasing $\langle M^2\rangle$.
Narrow disorder error bands do not exclude such a problem. The response enhancement must therefore
be considered together with the local maps and dynamical measurements below.

\clearpage
\subsection{Spatial structure of the local susceptibility}
The response of an individual spin to a uniform field is
\begin{equation}
\chi_{i,a}=\frac{\langle s_i M\rangle_a}{T}=\frac{1}{T}\sum_j\langle s_i s_j\rangle_a,\qquad \sum_i\chi_{i,a}=\chi_a.
\end{equation}
Figure~\ref{fig:local} follows the first $L=50$ disorder realization across four temperatures, so
the bond geometry remains fixed.
\begin{figure}[H]
\centering\includegraphics[width=0.76\linewidth]{02_local_susceptibility.png}
\caption{Local susceptibility estimates at $T=0.30,0.70,1.02,2.00$ for the same bond realization.
Each panel uses an independent symmetric linear color scale. Negative estimates are retained; equal
color intensity across panels does not mean equal response.}\label{fig:local}
\end{figure}
At $T=0.30$, broad patches have nearly uniform estimated response, indicating coherent spin behavior
over the measured trajectory. At $T=0.70$, these patches become more structured, while strong
spatial variation persists. At $T=1.02$, large positive responses are more concentrated. By
$T=2.00$, the scale has fallen from hundreds to values of order unity and finer variations dominate.
Residual heterogeneity is expected because local coordination and connectivity remain nonuniform
even when thermal correlations are short.

The negative patches require care. Ferromagnetic equilibrium correlations satisfy $\langle s_i
s_j\rangle_a\geq0$, implying $\chi_{i,a}\geq1/T$. Negative estimates cannot describe equilibrium
antiferromagnetic regions. If disconnected clusters rarely reverse during a finite trajectory, their
relative orientations do not average out; cross-cluster terms in $s_iM$ can remain spuriously
negative. This is a natural explanation for the coherent negative patches, although the stored map
does not identify every slow mode.

At $T=0.30$, approximately $44.5\%$ of sites have negative estimates and $40.8\%$ fall below minus
three times their estimated block SEM. These are not merely small fluctuations within the reported
errors. They indicate unresolved sampling or error-estimation problems. Thus the maps reveal spatial
structure in measured trajectories, but the low-temperature amplitudes cannot yet be treated as
converged equilibrium rare-region responses.

\clearpage
\subsection{Sample-to-sample susceptibility distributions}
Figure~\ref{fig:dist} shows kernel density estimates of $x=\log_{10}\chi$ from $1000$ realizations
at each size and temperature. Their normalization is
\begin{equation}
\int p_x(x)\,\mathrm dx=1,\qquad p_\chi(\chi)=\frac{p_x(\log_{10}\chi)}{\chi\ln10}.
\end{equation}
The ordinate is a density per unit logarithmic susceptibility, not per unit $\chi$.
\begin{figure}[H]
\centering\includegraphics[width=0.73\linewidth]{03_susceptibility_distributions.png}
\caption{Distributions of total susceptibility at six temperatures for $L=10,24,40,50$, normalized
in $\mathrm d(\log_{10}\chi)$. Smooth curves are kernel density estimates; small shoulders require
bandwidth and sampling checks before interpretation.}\label{fig:dist}
\end{figure}
Cooling moves the distributions toward larger susceptibility, consistent with Fig.~\ref{fig:mean}.
Their shapes reveal information hidden by the mean. At $T=0.70$, the larger sizes have particularly
broad distributions extending across several decades. Samples at identical $L$ and $T$ can therefore
respond very differently. Variations in cluster sizes and weak connections offer a plausible origin,
affecting both correlated magnetization and the time required to explore its orientations.

The width is not monotonic in temperature. The $T=0.30$ and $0.50$ curves concentrate at larger
responses and have appreciable tails toward smaller values, whereas the intermediate-temperature
distribution is especially broad. This may reflect a crossover between nearly rigid clusters and
thermally disrupted correlations. However, slow sampling can also reshape the distributions, and
these data do not separate the two contributions.

At $T=1.50$ and $2.00$, the distributions narrow as $L$ increases, consistent with averaging over
many weakly correlated regions. Their rightward shift is partly the trivial effect of total
susceptibility increasing with volume. Replacing $\chi$ by $\chi/N$ translates each logarithmic
distribution by $-\log_{10}N$ without changing its width.

A broad distribution can make the mean unrepresentative of typical samples, but does not establish a
power-law tail or a lack of self-averaging. Such conclusions require converged sampling and size
dependence of, for example, the relative variance. These densities contain both disorder variation
and MC estimation noise; their appearance alone cannot determine a Griffiths exponent.

\clearpage
\subsection{Autocorrelation under heat-bath dynamics}
A reference state is saved at $t_0=4096$ heat-bath sweeps after thermalization. The measured overlap
is
\begin{equation}
A_a(t)=\frac{1}{N}\sum_i s_i(t_0+t)s_i(t_0),\qquad 0\leq t\leq128.
\end{equation}
We use $A(t)$ to distinguish it from heat capacity; the original plot labels it $C(t)$. Each sample
contributes one fixed-origin trajectory, without time-origin averaging or subtraction of spin means.
\begin{figure}[H]
\centering\includegraphics[width=0.79\linewidth]{04_autocorrelation.png}
\caption{Disorder-averaged overlap with $\pm1$ SEM. One time unit is one random-site heat-bath
sweep. The reference state is taken $4096$ sweeps after thermalization, with $A(0)=1$. SW updates
are excluded from the time axis.}\label{fig:corr}
\end{figure}
Temperature has a much larger effect than the visible size dependence. At $L=50$, $[A(128)]_{\rm
dis}\simeq0.921,0.812,0.609,0.199$ for $T=0.30,0.50,0.70,1.02$, respectively. At $T=1.50$ and
$2.00$, the overlap is close to zero by the end of the window, with $T=2.00$ losing memory
substantially faster.

The low-temperature curves show a rapid initial drop followed by much slower decay. A cluster
picture suggests that weakly constrained spins relax first, while strongly correlated regions retain
memory. Reversing such regions through local updates may require rare intermediate configurations,
even though SW updates can flip their overall orientation efficiently. The plotted relaxation
therefore characterizes the chosen local heat-bath dynamics, not the autocorrelation time of an SW
sampling algorithm.

Similar curves across sizes indicate modest size dependence within this window, rather than
size-independent asymptotic relaxation. An apparent plateau over $128$ sweeps does not establish
frozen order or a nonzero infinite-time limit. For any finite system at positive temperature,
heat-bath dynamics can eventually sample both orientations. Longer trajectories are needed to
resolve the tail; the linear-axis plots do not distinguish power-law, stretched-exponential, or
other decay forms. The SEM combines disorder variation with single-trajectory noise, and adjacent
times are correlated.

\clearpage
\section{Summary}
The four figures provide complementary views of the $p=0.5$ bond-diluted Ising model. The heat
capacity identifies a broad energy-fluctuation crossover near $T=1$, with an approximately
size-independent per-site peak. Magnetic susceptibility is more sensitive to large correlated
clusters and grows strongly upon cooling. Its distributions show that this enhancement varies
substantially between disorder realizations, especially at intermediate temperatures.

The local maps and temporal overlap connect this behavior to spatial heterogeneity and slow
heat-bath dynamics. These observations are compatible with a cluster-based picture of the Griffiths
region, but at the percolation threshold critical connectivity must also be considered. Neither the
heat-capacity maximum nor finite-time memory identifies a finite-temperature ferromagnetic
transition.

A central limitation is that low-temperature local-response estimates violate equilibrium
positivity. Their magnitude relative to the reported block errors indicates that static sampling is
not uniformly converged. Narrow error bands on disorder averages do not remove this concern, and
apparent distribution tails may also be affected. The results therefore document finite-size
responses and finite-time relaxation, rather than establish a Griffiths singularity or a dynamical
exponent.

For quantitative conclusions, static observables should first be checked using efficient cluster
sampling and independent runs, while local heat-bath updates are retained for the specified
dynamics. Longer observation windows, multiple time origins or independent trajectories, and
block-length checks would then resolve slow relaxation more reliably. Separating equilibrium
sampling from the chosen stochastic dynamics is essential for identifying robust physical features.
\end{document}

